\documentclass[11pt]{article}
\usepackage[margin=1in]{geometry}
\usepackage{amsmath,amssymb,amsthm,mathtools}
\usepackage{booktabs}
\usepackage{enumitem}
\usepackage{xcolor}
\usepackage{hyperref}
\hypersetup{colorlinks=true, linkcolor=blue!50!black}
\newcommand{\E}{\mathbb{E}}
\newcommand{\Eqx}{\E_{q(x_i)}}
\newcommand{\Cov}{\operatorname{Cov}}
\newcommand{\tr}{\operatorname{tr}}
\newcommand{\code}[1]{\texttt{\small #1}}
\emergencystretch=3em \hbadness=10000
\title{The $q(x)$ update given $q(f)$: two evaluation paths}
\date{}
\begin{document}
\maketitle
\vspace{-2.5em}

\noindent
Picking up where \texttt{efgp\_estep\_succinct.tex} leaves off: the $q(f)$
step is done, so we hold the posterior-mean drift
$\bar f_r(x)=\phi(x)^*\mu_r$ fixed on the spectral grid $\{\xi_k\}_{k=1}^M$
(weights $D_k$, centre $x_c$), and take one SING natural-gradient step on
$q(x)$. Diffusion is isotropic, $\Sigma=\sigma^2 I$.

% ==========================================================================
\section{What we compute (common to both paths)}
% ==========================================================================
\begin{enumerate}[leftmargin=1.6em,itemsep=2pt]
\item By exponential-family duality the natural gradient is
  $\partial\mathcal L/\partial\mu$, $\mu=(\E[x_i],\E[x_ix_i^\top],\E[x_ix_{i+1}^\top])$.
\item The transition term depends on the drift only through three
  per-source moments under $q(x_i)=\mathcal N(m_i,S_i)$:
  \[
  \mathrm{Ef}_i=\Eqx[\bar f],\quad
  \mathrm{Edfdx}_i=\Eqx[J_{\bar f}],\quad
  \mathrm{Eff}_i=\Eqx[\bar f^\top\Sigma^{-1}\bar f].
  \]
\item \textbf{Method.} The three moments are \emph{per source}
  ($\mathrm{Ef}_i,\mathrm{Edfdx}_i,\mathrm{Eff}_i$ for every $i$). From each
  we assemble the per-transition scalar $\code{trm}_i$
  (\code{compute\_neg\_CE\_single}) using the $q(x)$ statistics
  $(d_i,C_i,m_i,S_i,S_{i,i+1})$, then sum: $\mathcal L(\mu)=\sum_i\code{trm}_i$.
  The sum is ours --- formed only to hand \code{jax.grad} a scalar; since
  $\code{trm}_i$ depends only on $\mu_i$, \code{jax.grad}$(\mathcal L)$ once
  returns every per-transition $(J,h,L)$. No hand-derived gradients and
  nothing dropped --- reverse mode returns the exact gradient.
\item Each moment is a Gaussian average of Fourier features, i.e.\ a
  characteristic function. With the enveloped feature
  \[
  \psi_k(m,S)\;=\;D_k\,e^{2\pi i\,\xi_k^\top(m-x_c)}\;e^{-2\pi^2\,\xi_k^\top S\,\xi_k},
  \]
  \[
  \Eqx[\bar f_r]=\sum_k \mu_{r,k}\,\psi_k,\qquad
  \Eqx[\partial_j\bar f_r]=\sum_k \mu_{r,k}\,(2\pi i\,\xi_{k,j})\,\psi_k,
  \]
  \[
  \Eqx[\bar f^\top\bar f]=\sum_\delta \rho(\delta)\,
    e^{2\pi i\,\delta^\top(m-x_c)}\,e^{-2\pi^2\,\delta^\top S\,\delta},
  \quad \rho(\delta)=\sum_r\!\!\sum_{\xi_k-\xi_{k'}=\delta}\!\! (D_k\mu_{r,k})(D_{k'}\mu_{r,k'}).
  \]
\end{enumerate}
The two paths differ \emph{only} in how these sums --- and, under
\code{jax.grad}, their adjoints --- are evaluated. $N=K(T{-}1)$ is the
source count.

% ==========================================================================
\section{Path A --- direct sums \ ($O(NM)$)}
% ==========================================================================
\begin{enumerate}[leftmargin=1.6em,itemsep=2pt]
\item Increments $d_i=m_{i+1}-m_i$, $C_i=S_{i,i+1}-S_i$.
\item Autocorrelation $\rho(\delta)$: one padded FFT of $\{D_k\mu_{r,k}\}$
  per output $r$, summed over $r$ (once per outer iter).
\item Moments, \emph{one per source} $i$, by direct summation:
  $\mathrm{Ef}_i,\mathrm{Edfdx}_i$ as the $O(M)$ spectral sums above,
  $\mathrm{Eff}_i$ as the $O(M)$ sum over the difference grid against
  $\rho$ --- each a value at $m_i$, no sum over sources. Vectorised over
  all $N$ sources $\Rightarrow O(NM)$.
\item Assemble $\mathcal L=\sum_i\code{trm}_i$; \code{jax.grad}$(\mathcal L)$
  w.r.t.\ $\mu$. Reverse mode differentiates the dense sums directly
  (another $O(NM)$ pass); no NUFFT.
\end{enumerate}
\textbf{Cost} $O(NM)$ time and memory (the reverse tape holds the
$(N,M)$ mode arrays). \textbf{Exact.} Simplest; best at small $M$.
Implementation: \code{exp\_full\_moments.py} + \code{exp\_batched\_estep.py}.

% ==========================================================================
\section{Path B --- gather / NUFFT \ ($O(M\log M + N)$)}
% ==========================================================================
\begin{enumerate}[leftmargin=1.6em,itemsep=2pt]
\item Increments $d_i,C_i$ as above.
\item Coefficient grids: $c^{(0)}_k=D_k\mu_{r,k}$ (mean),
  $c^{(1)}_{k,j}=(2\pi i\,\xi_{k,j})D_k\mu_{r,k}$ (Jacobian); and the
  difference-grid $\rho(\delta)$ (one padded-FFT autocorrelation) for the
  quadratic.
\item Evaluate the moments by the \emph{gmix gather} --- a
  Gaussian-enveloped Type-2 NUFFT that reads the coefficient field at each
  $m_i$ with its own $\mathcal N(\,\cdot\,;m_i,S_i)$ stencil, returning
  \emph{one value per source} (not a sum): one batched call for
  $\mathrm{Ef}$, one per latent dim for $\mathrm{Edfdx}$, one on the
  $\rho$-grid for $\mathrm{Eff}$. Cost $O(M\log M + N)$ per call. (Contrast
  the $q(f)$-side \emph{spreader}, a Type-1 NUFFT that \emph{accumulates}
  $\sum_i w_i\,\mathcal N(\cdot;m_i,S_i)$ into the BTTB generator --- that
  is the primitive that sums over sources; the gather does not.)
\item Assemble the per-source $\code{trm}_i$ and $\mathcal L=\sum_i\code{trm}_i$;
  \code{jax.grad}$(\mathcal L)$. The VJP of each gather is its Type-1 adjoint
  --- the \emph{spreader} (\code{efgp\_gmix\_spreader}), the same
  $\sum_i$-onto-the-grid primitive used on the $q(f)$ side --- chained
  automatically, also $O(M\log M + N)$, memory flat in $M$.
  So the whole per-source pipeline (moments $\to \code{trm}_i \to$ gradients)
  is $O(M\log M + N)$: the gather (Type-2) and spreader (Type-1) are an
  adjoint pair, forward and backward.
\end{enumerate}
\textbf{Cost} $O(M\log M + N)$ time, memory independent of $M$.
\textbf{Approximate}: the per-source Gaussian stencil truncates at
$n_\sigma$ (tunable \code{gather\_N}, \code{stencil\_r}); error decreases
on finer grids. Best at large $M$/$N$. Primitive:
\code{efgp\_gmix\_gather.py} / \code{drift\_moments\_gmix\_jax}.

% ==========================================================================
\section{Choosing a path}
% ==========================================================================
\begin{center}
\renewcommand{\arraystretch}{1.25}
\begin{tabular}{@{}lll@{}}
\toprule
 & Path A (direct) & Path B (gather) \\
\midrule
time    & $O(NM)$          & $O(M\log M + N)$ \\
memory  & $O(NM)$          & flat in $M$ \\
moments & exact            & stencil-approx ($n_\sigma$) \\
grad    & autodiff (dense) & autodiff (NUFFT adjoints) \\
best when & $M$ small       & $M$ or $N$ large \\
\bottomrule
\end{tabular}
\end{center}
\noindent
Both return the same natural gradient (to the gather's stencil tolerance).
The $O(NM)$ cost is exclusive to Path A (brute-force sums, no gmix); Path B's
gather/spreader pair keeps \emph{per-source} moments and gradients at
$O(M\log M + N)$.
Empirically the crossover is $M\!\approx\!200$--$350$: below it A is faster
and exact; above it B is several-fold faster with flat memory. In neither
path is any gradient hand-derived or dropped --- both are \code{jax.grad}
of the exact summed transition term; they differ only in the primitive
used to evaluate (and hence adjoint) the moment sums.
\end{document}
